% 31 Aug

Note that in what follows, when we say ``$\sup$'', we mean ``union over''. Just like in Lean.

\begin{boxproposition}
    $\forall y \parenth{y \in \WF \lr y \subset \WF}$
\end{boxproposition}
\begin{proof}\hfill
    \begin{description}
        \item[$\to$] Let $\beta = \rank{y}$. Then $y \ssq V_{\beta} \subset \WF$.
        \item[$\leftarrow$] Assume $y \subset \WF$. Define $f : y \to \OR$ by $f(x) = \rank{x} + 1$. We're using replacement to know that $f \in V$.

        Let $\beta = \sup_{x \in y}{f(x)}$. Then, $y \ssq V_{\beta}$. So $y \in \Powset{V_{\beta}} = V_{\beta + 1} \subset \WF$.
    \end{description}
\end{proof}

\begin{boxproposition}
    \begin{align*}
        \forall y \in \WF,\, \forall x \in y,\, \brac{x \in \WF \land \rank{x} < \rank{y}}
    \end{align*}
\end{boxproposition}
\begin{proof}
    We know $x \in y \ssq V_{\rank{y}} \subset \WF$, so $x \in \WF$. This also shows that $\rank{x} + 1 \leq \rank{y}$, so $\rank{x} < \rank{y}$.
\end{proof}

\begin{boxproposition}[The Rank Formula]
    \begin{align*}
        \forall y \in \WF, \, \rank{y} = \sup_{x \in y}\of{\rank{x} + 1}
    \end{align*}
\end{boxproposition}
\begin{proof}
    Put $\rho = \rank{y}$ and $\sigma = \sup_{x \in y}\of{\rank{x} + 1}$. By previous prop, if $f \in y$, then $\rank{x} < \rank{y} = \rho$, so $\rank{x} + 1 \leq \rho$. So $\sigma \leq \rho$. % [CLAUDE] add cross-reference here

    It suffices to show that $y \ssq V_{\sigma}$ to conclude that $\rho \leq \sigma$. Now $V_{\sigma} = \setst{x}{\rank{x} < \sigma}$. So $y$ is indeed a subset. % [CLAUDE] why? elaborate
\end{proof}

\begin{boxproposition}
    \begin{align*}
        \forall \beta \in \OR, \,
        \brac{\beta \in \WF \land \rank{\beta} = \beta}
    \end{align*}
\end{boxproposition}
\begin{proof}
    We argue by induction on $\beta$. Base case is trivial. In the inductive case, $\beta \ssq \WF$, so by previous prop, $\beta \in \WF$. So,
    \begin{align*}
        \rank{\beta}
        &= \sup_{\alpha < \beta}\of{\rank{\alpha} + 1} \\
        &= \sup_{\alpha < \beta}\of{\alpha + 1} \\
        &= \beta
    \end{align*}
\end{proof}

\begin{boxproposition}
    \begin{align*}
        \forall \beta \in \OR, \,
        \brac{V_{\beta} \cap \OR = \beta}
    \end{align*}
\end{boxproposition}
\begin{proof}
    We know
    \begin{align*}
        V_{\beta} \cap \OR
        &= \setst{x \in \WF}{\rank{x} < \beta} \cap \OR \\
        &= \setst{\alpha \in \OR}{\rank{\alpha} < \beta} \\
        &= \setst{\alpha \in \OR}{\alpha < \beta} \\
        &= \beta
    \end{align*}
\end{proof}

\begin{boxproposition}
    \begin{align*}
        \forall n < \omega, \,
        V_{n} \text{ is finite}
    \end{align*} % [CLAUDE] maybe reformat this a little
\end{boxproposition}

Indeed, their sizes are $0$, $2^0$, $2^{2^{0}}$, and so on.

\begin{boxproposition}
    $V_{\omega}$ is countable.
\end{boxproposition}
\begin{proof}[Proof using AC]
    Write
    \begin{align*}
        V_{\omega} = \bigcup_{n < \omega} V_n
    \end{align*}
    AC tells us that a countable union of countable sets is countable.
\end{proof}
\begin{proof}[Proof that avoids AC]
    By recursion on $n < \omega$, we can define $R_n \ssq V_n \times V_n$ so that
    \begin{enumerate}
        \item $\parenth{V_n, R_n}$ is a well-ordering
        \item $R_{n+1}$ end-extends $R_n$
    \end{enumerate}
    Then, we'll be able to put $R_{\omega} = \bigcup_{n < \omega} R_n$ and show taht $R_{\omega}$ is a well-ordering of type $\omega$.

    For $x, y \in V_{n + 1}$, put $x \, R_{n+1} \, y$ iff
    \begin{enumerate}
        \item $x, y \in V_n$ and $x \, R_n \,, y$.
        \item $x \in V_n$ and $y \in V_{n + 1} \setminus V_n$.
        \item $x, y \in V_{n+1} \setminus V_n$ 
    \end{enumerate}
    \sorry
    % [CLAUDE] Here's a hint. You can complete this by essentially defining $R_{n+1}$ as some lexicographic ordering or something, I mean write $x$ as a bunch of (lexicographically) related things and then... something, not really sure...
\end{proof}

\begin{remark}
    If ZF is consistent, then so is ZF $+$ ``there's a countable union of pairs that has no cardinality''.
\end{remark}

\begin{remark}
    Without AC, we don't know whether the set $V_{\omega + 1} = \Powset\of{V_{\omega}}$ even has a \textit{cardinality}. With AC,
    \begin{align*}
        \abs{V_{\omega + 1}} = \abs{\Powset\of{V_{omega}}} = 2^{\aleph_0}
    \end{align*}
    where $\aleph_0 = \omega = \abs{V_{\omega}}$.
\end{remark}

% [CLAUDE] in what follows, replace k with \kappa
Given AC, we know that $2^k > k$ and we can define $k^{+}$ to be the least cardinali strictly greater than $k$. Without AC, we still know, given a cardinal $k$, that
\begin{align*}
    \setst{R \ssq k}{(k, R) \text{ is a well-ordering}}
\end{align*}
is a set. We can still have Mostowski Collapse for well-orderings. We can then conclude that
\begin{align*}
    \setst{\text{types}\of{k, R}}{(k, R) \text{ is a well-ordering}}
\end{align*}
is a set of ordinals. Then, we can set $k^{+}$ to simply be the supremum of this set.

Without AC, define
\begin{align*}
    \aleph_0 &:= \omega \\
    \aleph_{\alpha + 1} &:= \parenth{\aleph_{\alpha}}^{+} \\
    \aleph_{\beta} &:= \sup_{\alpha < \beta}\of{\aleph_{\alpha}} \text{ if } \beta \text{ is a limit ordinal} 
\end{align*}
With AC, define
\begin{align*}
    \beth_0 &:= \omega \\
    \beth_{\alpha + 1} &:= 2^{\beth_0} \\
    \beth_{\beta} &:= \sup_{\alpha < \beta}\of{\beth_{\alpha}} \text{ if } \beta \text{ is a limit ordinal}
\end{align*}

It's clear 